\section{Глава~\ref{sec:neurontypes}. Типы нейронов}

Числа главы держатся на прямых подсчётах клеток в мозге человека там, где такие подсчёты есть, правило «один нейрон --- один медиатор» --- на клеточных атласах и опытах, которые нашли и исключения, роль \nt{ДОФА} --- на записях импульсов, а роли остальных широковещательных каналов --- на гипотезах.

{\footnotesize
\setlength{\tabcolsep}{3pt}\renewcommand{\arraystretch}{1.15}
\begin{longtable}{>{\raggedright}p{0.5cm}>{\raggedright}p{4.0cm}>{\raggedright}p{5.6cm}>{\raggedright}p{2.3cm}>{\raggedright\arraybackslash}p{1.7cm}}
\toprule
№ & Утверждение & Опыт и что измерено & Источник & Статус \\
\midrule\endfirsthead
\toprule
№ & Утверждение & Опыт и что измерено & Источник & Статус \\
\midrule\endhead
1 & В мозге человека около $8.6\cdot10^{10}$ нейронов, и почти все \nt{ГЛУТ}-нейроны --- мелкие нейроны мозжечка (табл.~\ref{tab:types}) & Изотропный фракционатор: ткань растворяют в однородную взвесь ядер и считают ядра с нейронным маркером NeuN. У взрослых мужчин $86.1\pm8.1$ млрд нейронов; в коре больших полушарий из них лишь $19\,\%$, основная масса --- в мозжечке & \cite{Azevedo2009} & измерено \\
2 & У почти каждого нейрона один главный медиатор (утверждение~\ref{prop:dale}), но есть исключения & Транскриптомный атлас мозга мыши: около $4\cdot10^6$ клеток, $5322$ кластера; кластерам приписан медиатор по генам синтеза и упаковки. Исключения измерены прямо: \nt{ДОФА}-нейроны в срезах выбрасывают ещё и ГАМК через тот же везикулярный транспортёр и тормозят нейроны стриатума; у части \nt{ДОФА}-аксонов глутамат и дофамин выходят из разных окончаний одного провода (электронная микроскопия и оптогенетика) & \cite{Strata1999,Yao2023,Tritsch2012,Zhang2015} & измерено \\
3 & Весь столбец матрицы весов одного знака~\eqref{eq:dalesign} & Вывод в главе из утверждения~\ref{prop:dale} и того, что рецепторы глутамата почти все возбуждающие, а ГАМК --- тормозные. Прямой проверки «все получатели одного нейрона получают вес одного знака» как общего правила нет; контрпримеры --- совместный выброс ГАМК \nt{ДОФА}-нейронами (строка~2) и получатели с рецепторами разного знака (строка~11) & вывод в главе & теория \\
4 & От трёх четвертей до четырёх пятых нейронов коры --- \nt{ГЛУТ}, от пятой части до четверти --- \nt{ГАМК} & Иммуноокраска на ГАМК в столбиках шириной $50$\,мкм через всю толщу коры в $10$ областях коры обезьяны: в $8$ областях ГАМК-нейроны составляют около $25\,\%$ всех нейронов, в первичной зрительной коре --- меньше $20\,\%$ & \cite{Hendry1987} & измерено \\
5 & У \nt{ГЛУТ}-нейрона около $10^4$ выходов & Стереология на вскрытии: в новой коре молодых мужчин $1.64\cdot10^{14}$ синапсов (электронная микроскопия, дисектор) и около $2.3\cdot10^{10}$ нейронов. Отношение даёт $\sim7\cdot10^3$ синапсов на нейрон; в среднем число входов равно числу выходов & \cite{Tang2001synapses,Pakkenberg1997} & измерено \\
6 & Выход \nt{ДОФА}-нейрона ветвится на $10^5$--$10^6$ окончаний; это оценка по охвату мишени, а не подсчёт & Вирусная метка мембраны отдельных \nt{ДОФА}-нейронов крысы и полная реконструкция аксона: один нейрон покрывает от $0.45$ до $5.7\,\%$ (в среднем $2.7\,\%$) объёма стриатума. Измерен охват, а не число окончаний: оно выведено из охвата по порядку величины, прямого счёта окончаний нет. Для \nt{СЕРО} и \nt{НОРА} числа окончаний --- грубые оценки & \cite{Matsuda2009} & измерено \\
7 & Широковещательных нейронов мало: \nt{ДОФА} $\sim5\cdot10^5$ (главная из двух групп, нижняя граница), \nt{СЕРО} $\sim3\cdot10^5$ (сумма двух частичных подсчётов), \nt{ГИСТ} $\sim6\cdot10^4$ (табл.~\ref{tab:types}, рис.~\ref{fig:typepie}) & Стереологические подсчёты у человека: $550\,000$ пигментированных нейронов в чёрной субстанции (без вентральной покрышки); $235\,000$ нейронов в дорсальном ядре шва (все нейроны ядра) и ещё $125\,000$ серотониновых нейронов в покрышке моста; около $64\,000$ гистаминовых нейронов гипоталамуса. Для \nt{НОРА}, \nt{ГЛИЦ}, \nt{АЦЕТ} и \nt{АДРЕ} прямого подсчёта нет: в таблице это грубые оценки по порядку величины & \cite{Pakkenberg1991,Baker1990,Baker1991,Airaksinen1991} & измерено \\
8 & Глутамат в щели взлетает примерно до миллимоля и спадает за $\sim1\ms$ & По вытеснению быстро отсоединяющегося конкурентного блокатора с рецепторов NMDA во время синаптической передачи (культура нейронов гиппокампа): пик $1.1$\,мМ, постоянная спада $1.2\ms$ & \cite{Clements1992} & измерено \\
9 & В работающей коре возбуждающий и тормозящий токи почти уравновешены, нейрон сидит у порога & Теория: сеть с сильными связями сама приходит в уравновешенный режим. Опыт: в префронтальной коре хорька in vivo во время «вверх»-состояний потенциал реверсии синаптического тока держится около $-37$\,мВ сотни миллисекунд, хотя проводимость меняется на $\sim21$\,нСм: возбуждение и торможение растут и падают вместе & \cite{vanVreeswijk1996,Haider2006} & измерено \\
10 & \nt{ДОФА}-нейроны сообщают ошибку предсказания награды~\eqref{eq:rpe} (рис.~\ref{fig:rpe}) & Импульсы \nt{ДОФА}-нейронов обезьяны: пачка на неожиданный сок; после обучения --- пачка на сигнал и нет ответа на предсказанный сок; провал частоты, когда обещанный сок не дан. В количественном опыте частота пропорциональна разности текущей награды и взвешенного среднего прошлых, но только для положительных ошибок. Само уравнение~\eqref{eq:rpe} --- алгоритм временных разностей & \cite{Schultz1997,Bayer2005,SuttonBarto2018} & измерено \\
11 & Знак и скорость задаёт рецептор: ионотропные --- доли миллисекунды, метаботропные --- гораздо медленнее (утверждение~\ref{prop:receptor}) & Быстрые импульсы глутамата на кусочки мембраны нейронов гиппокампа: ток через ионотропные рецепторы нарастает (от $20$ до $80\,\%$) за $0.2$--$0.6\ms$ и спадает за $2$--$3\ms$. Медленный тормозный потенциал пирамидных нейронов снимается избирательным блокатором метаботропного рецептора ГАМК. Два семейства рецепторов дофамина с противоположным действием; в стриатуме нейронам с рецептором D1 дофамин нужен для усиления синапсов, нейронам с D2 --- для ослабления & \cite{Colquhoun1992,DutarNicoll1988,Beaulieu2011,Shen2008} & измерено \\
12 & Широковещательный канал --- не один провод, а жгут из немногих линий (раздел~\ref{sec:buses}) & Вирусно-генетическая разметка серотониновых нейронов дорсального ядра шва мыши: нейроны, посылающие выход в миндалину и в лобную кору, лежат в разных местах ядра, ветвятся в дополняющие друг друга области и противоположно отвечают на неприятный раздражитель. Обзор: подгруппы нейронов голубого пятна (\nt{НОРА}) кодируют разные процессы & \cite{Ren2018,Poe2020} & измерено \\
13 & \nt{НОРА} задаёт коэффициент усиления $g$ & Гипотеза адаптивного усиления; модель сети, в которой катехоламины повышают крутизну передаточной характеристики. Прямого измерения «$g$ растёт с норадреналином» во всей сети нет; измерено, что диаметр зрачка следует за импульсами нейронов голубого пятна обезьяны & \cite{AstonJones2005,ServanSchreiber1990,Joshi2016} & гипотеза \\
14 & \nt{АЦЕТ} переключает «запись/чтение» и задаёт скорость обучения & Гипотеза. Опора: в срезах обонятельной коры крысы агонисты ацетилхолина ($100$\,мкМ) ослабляют внутренние, «вспоминающие» синапсы больше чем на $60\,\%$, а входные --- меньше чем на $15\,\%$ & \cite{Hasselmo2006,HasselmoBower1992} & гипотеза \\
15 & \nt{СЕРО} задаёт коэффициент дисконтирования $\gamma$; серотониновые нейроны работают ровно, несколько импульсов в секунду, и почти молчат в одной из фаз сна & Гипотеза «серотонин $=\gamma$». Сильнейший довод: оптогенетическая активация серотониновых нейронов дорсального ядра шва мыши уменьшает число преждевременных отказов от ожидания отложенной награды. Частота импульсов и молчание во сне с быстрыми движениями глаз измерены (обзор записей) & \cite{Doya2002,Miyazaki2014,JacobsAzmitia1992} & гипотеза \\
\bottomrule
\end{longtable}}
